The servants have found something genuinely wrong in the field. The owner
sowed good seed, yet another crop now grows among it. They offer to put the
matter right by pulling out the weeds. His refusal concerns the wheat:
their attempt to save it could destroy it. He orders them to wait until
harvest, when he will command a separation. The Gospel therefore presents
both a prohibition of premature uprooting and a promise of final judgement.
The same owner gives both commands.

Paul describes the life Christians must lead during this interval. The
Epistle asks people with grievances against one another to forgive as the
Lord has forgiven them. It also commands them to teach and admonish one
another. Charity, correction and patience belong to the same community.
The questions are practical as well as exegetical: what distinguishes
faithful care for another person from the servants' destructive haste,
and how can people who still need mercy exercise that care?

The Mass's prayers repeatedly acknowledge this need. The Collect asks for
continuing protection. The Secret asks pardon and the guidance of wavering
hearts. After Communion, those who have received salvation's pledge still
ask for its effect. Divine promises have not made petition unnecessary.
The Introit already places God's promise of peace beside the people's
calling upon him, and the Communion appoints Christ's assurance that
believing prayer will receive.

Two interpretations follow the detailed commentary. The first develops
Chrysostom's reasons for sparing the tares and Augustine's warning that
human beings can mistake wheat for weeds. It considers forgiveness and
correction within a community whose members can still change. The second
begins with Jerome's question about Jeremias: why does God command prayer
for a restoration he has already promised? Augustine's account of the
sinner crying from the depths explains how such confidence depends on
Christ's mercy. These interpretations meet in conversion and hope, while
answering different difficulties: the first concerns our action toward
another person; the second concerns our asking and receiving from God.

This is the resumed Fifth Sunday after the Epiphany in the 1962 Roman
Missal, appointed on 8 November 2026 in the General Roman Calendar.
The printed Mass is on pp.~420--421, with the Preface of the Trinity
on p.~293. There are eleven proper elements and no optional branches
or additional commemorative orations in this general-calendar occurrence.

\section{The propers at a glance}
\begin{maptable}
Introit & Jer 29:11,12,14; Ps 84:2 (Hebrew 85) & God promises peace, hearing and regathering; the psalm recalls restoration.\\
Collect & \textit{Familiam tuam} & God's family asks continual protection, relying on heavenly grace alone.\\
Epistle & Col 3:12--17 & Forgiveness, charity, one-body peace, instruction, song and thanksgiving.\\
Gradual & Ps 43:8--9 (Hebrew 44) & Rescue is ascribed to God and answered with enduring praise.\\
Alleluia & Ps 129:1--2 (Hebrew 130) & The cry from the depths asks the Lord to hear.\\
Gospel & Matt 13:24--30 & Wheat and tares grow together until the owner's harvest.\\
Offertory & Ps 129:1--2 (Hebrew 130) & The opening cry from the depths is repeated.\\
Secret & \textit{Hostias tibi} & The offerers ask pardon and direction for wavering hearts.\\
Preface & Of the Trinity & Thanksgiving confesses distinct Persons, one essence and equal majesty.\\
Communion & Mark 11:24 & Christ promises reception to believing prayer.\\
Postcommunion & \textit{Quaesumus, omnipotens Deus} & Those who received the pledge ask for salvation's effect.\\
\end{maptable}
\tableofcontents
